\documentclass[main.tex]{subfiles}
\begin{document}

\subsection{Problem 1}
Sketch the Chisel code for the circuit in the figure. The register shall be 8 bits wide and shall be initialized to 0 on reset.

Explain in your own words what the circuit does.
\begin{figure} [h]
    \centering
    \includegraphics[width=0.5\linewidth]{images//2022/Q1-2022.png}
\end{figure}

\subsubsection{Solution}
\begin{adjustwidth}{2em}{0pt}
    \begin{multicols}{2}
    \begin{verbatim}
        import chisel3._
        import chisel3.util._
        
        class Circuit extends Module {
            val io = IO(new Bundle {
                val sel  = Input(UInt(2.W))
                val regAcc = Output(UInt(8.W))
            })
        
            val regAcc = RegInit(0.U(8.W))

            // Adders
            val plus1 = regAcc + 1.U
            val plus3 = regAcc + 3.U

            val muxVal = MuxLookup(io.sel, regAcc)(
                Seq(
                    0.U -> regAcc,
                    1.U -> 0.U,
                    2.U -> plus1,
                    3.U -> plus3
                )
            )

            regAcc := muxVal
            io.regAcc := regAcc
        }
    \end{verbatim}
    \columnbreak
    \begin{verbatim}
        import chisel3._
        import chisel3.util._

        class Circuit extends Module {
            val io = IO(new Bundle {
                val sel  = Input(UInt(2.W))
                val regAcc = Output(UInt(8.W))
            })

            val regAcc = RegInit(0.U(8.W))

            switch(sel) {
                is(0.U) {regAcc := regAcc}
                is(1.U) {regAcc := 0}
                is(2.U) {RegAcc := regAcc + 1.U}
                is(3.U) {RegAcc := regAcc + 3.U}
            }
            
            io.regAcc := regAcc
    \end{verbatim}  
    \end{multicols}
    The system consists of an register than through the select signal Sel, can be set to 0, or increment by +1, or +3. 

    The output is the stored value of the register.
\end{adjustwidth}

\newpage
\subsection{Problem 2}
Given following Chisel code:
\begin{verbatim}
    class X extends Module {
        val io = IO(new Bundle{
            val o = Output(UInt(32.W))
        })
        val r = RegInit(0.U(32.W))
        val s = RegInit(0.U(32.W))
        r := r + 1.U
        when (r === 100000.U) {
            r := 0.U
            s := s + 1.U
        }
        io.o := s
    }
\end{verbatim}
Explain the function and draw a block diagram.

\subsubsection{solution}
\begin{figure}[h]
    \centering
    \includegraphics[width=0.5\linewidth]{images//2022/Q2-2022.png}
\end{figure}
\begin{adjustwidth}{2em}{0pt}
    This system consists of two registers, r and s. The r-reister starts at 0 and counts to 10000 before reseating to 0.
    
    When the r-register resets to 0, the s-register increments by 1, the output is the value stored by the s-register.
\end{adjustwidth}

\newpage
\subsection{Problem 3}
\begin{multicols}{2}
    \includegraphics[width=1\linewidth]{images//2022/Q3-circuit-2022.png}
    \columnbreak
    \includegraphics[width=1\linewidth]{images//2022/Q3-timing-2022.png}
\end{multicols}
Find and describe the critical path in the circuit shown in the figure.
What is the maximum clock frequency?

\subsubsection{Solution}
\begin{adjustwidth}{2em}{0pt}
    There are 4 routes, that are equaly long:\\
    R1-G2-G4-G5-R5\\
    R1-G2-G3-G7-R6\\
    R2-G2-G4-G5-R5\\
    Rr-G2-G3-G7-R6

    Calculating the time we get, DFF clock-to-Q delay + 4 gate delays + setup time. With tCQ prop = 11ps, gate prop = 7ps, tSU = 3ps.
    \begin{align}
        11\text{ ps} + 3 \cdot 7\text{ ps} + 3\text{ ps} = 35\text{ ps}
    \end{align}

    From this we can find the maximum clock frequency with: 
    \begin{align}
        f_{\max} = \frac{1}{T_{\min}} = \frac{1}{35 \cdot 10^{-12}\text{ s}} = 28.6\text{ GHz}
    \end{align}
\end{adjustwidth}

\newpage
\subsection{Problem 4}
Implement a timer. It has one input: start and one output: done. The timer is started
with a single clock cycle high (1) of input start. It shall count for 10 clock cycles and then
output high (1) at output done. During counting done shall be low (0).
Draw the schematic of the circuit and sketch the Chisel code to describe it.

\subsubsection{Solution}
\begin{figure} [h]
    \centering
    \includegraphics[width=0.5\linewidth]{images//2022/Q4-2022.png}
\end{figure}
\begin{adjustwidth}{2em}{0pt}
    \begin{verbatim}
        import chisel3._
        import chisel3.util._
        
        class Circuit extends Module {
            val io = IO(new Bundle {
                val start  = Input(Bool())
                val done = Output(Bool())
            })
        
            val cnt = RegInit(0.U(4.W))

            val done = cnt === 0.U
            val dec = cnt - 1.U

            cnt := Mux(io.start,
                    10.U, 
                    Mux(done,
                    0.U,
                    dec
                )
            )
            
            io.done := done
        }
    \end{verbatim}
\end{adjustwidth}

\newpage
\subsection{Problem 5}
Draw the state machine and the datapath for the following circuit:
\begin{itemize}
    \item A vending machine that accepts coins of 5 and 10 kr. A soda costs 10 kr.
    \item When the sum of coins inserted into the vending machine is larger or equal to 10 kr, one can buy a soda by pressing the buy button.
    \item The inputs to the circuit are coin5, coin10, buy, and return.
    \item By pressing the return button, any money in the machine will be returned to the customer.
    \item Assume that all input signals are valid for one clock cycle (no edged detection needed).
    \item The outputs are a release and a money Return signal. The output signals have to be valid for only one clock cycle.
\end{itemize}
Note, that this exercise is related to your vending machine project, but it is NOT your
vending machine. The specification is different (considerably simpler).

\subsubsection{Solution}
\begin{figure} [h]
    \centering
    \includegraphics[width=0.5\linewidth]{images//2021/Q4-state-2021.png}
\end{figure}
\begin{figure} [h]
    \centering
    \includegraphics[width=0.7\linewidth]{images/2021/Q4-datapath-2021.png}
\end{figure}

\end{document}
